\documentclass[12pt,a4paper,spanish]{report}
\usepackage[utf8]{inputenc}
\usepackage[spanish]{babel}
\usepackage{geometry}
\usepackage{graphicx}
\usepackage{booktabs}
\usepackage{longtable}
\usepackage{listings}
\usepackage{xcolor}
\usepackage{hyperref}
\usepackage{fontawesome5}
\usepackage{titlesec}
\usepackage{fancyhdr}
\usepackage{tocloft}
\usepackage{enumitem}
\usepackage{caption}
\usepackage{subcaption}
\usepackage{float}
\usepackage{csquotes}
\usepackage{datetime2}

% Configuración de página
\geometry{margin=2.5cm, headheight=15pt}
\setlength{\parskip}{0.5em}
\setlength{\parindent}{0pt}

% Colores corporativos
\definecolor{celestedep}{HTML}{11698E}
\definecolor{celestemid}{HTML}{3AA8C9}
\definecolor{celestelight}{HTML}{6FC1FF}
\definecolor{crem}{HTML}{FFF8E7}
\definecolor{yellowmid}{HTML}{FFE066}
\definecolor{textdark}{HTML}{1A1A2E}
\definecolor{textmuted}{HTML}{6B7280}

% Configuración de hipervínculos
\hypersetup{
    colorlinks=true,
    linkcolor=celestedep,
    filecolor=celestemid,
    urlcolor=celestemid,
    citecolor=celestedep,
    pdftitle={Documentación Técnica - Landing María Fernanda},
    pdfauthor={Henry Paolo Alfaro Sotil},
    pdfsubject={Documentación completa del proyecto Landing María Fernanda - Psicología Clínica ACT},
    pdfkeywords={Firebase, Firestore, Quill.js, Landing Page, Psicología, ACT, Web Development}
}

% Estilo de capítulos y secciones
\titleformat{\chapter}[display]
    {\normalfont\huge\bfseries\color{celestedep}}
    {\chaptertitlename\ \thechapter}
    {20pt}
    {\titlerule\vspace{10pt}}
    [\vspace{10pt}\titlerule]

\titleformat{\section}
    {\normalfont\Large\bfseries\color{celestedep}}
    {\thesection}
    {1em}
    {}

\titleformat{\subsection}
    {\normalfont\large\bfseries\color{celestemid}}
    {\thesubsection}
    {1em}
    {}

% Estilo de lista de código
\lstdefinestyle{codeStyle}{
    backgroundcolor=\color{crem},
    basicstyle=\ttfamily\small,
    breaklines=true,
    frame=single,
    framesep=5pt,
    rulecolor=\color{celestelight},
    commentstyle=\color{textmuted}\itshape,
    keywordstyle=\bfseries\color{celestedep},
    stringstyle=\color{yellowmid},
    numbers=left,
    numberstyle=\tiny\color{textmuted},
    stepnumber=1,
    numbersep=8pt,
    showspaces=false,
    showstringspaces=false,
    tabsize=4
}

\lstset{style=codeStyle}

% Header y footer
\pagestyle{fancy}
\fancyhf{}
\fancyhead[L]{\leftmark}
\fancyhead[R]{\thepage}
\fancyfoot[C]{\textit{Landing María Fernanda — Documentación Técnica v1.0}}
\renewcommand{\headrulewidth}{0.5pt}
\renewcommand{\footrulewidth}{0.5pt}
\renewcommand{\chaptermark}[1]{\markboth{#1}{}}

% Portada
\begin{document}

\begin{titlepage}
    \centering
    \vspace*{2cm}

    {\Huge\bfseries\color{celestedep} Landing María Fernanda\par}
    \vspace{0.5cm}
    {\Large\color{celestemid} Psicología Clínica ACT — Documentación Técnica Completa\par}

    \vspace{2cm}

    {\large\normalfont Documentación para cliente y desarrolladores\par}

    \vspace{3cm}

    \begin{minipage}{0.8\textwidth}
        \centering
        \begin{tabular}{ll}
            \textbf{Autor:} & Henry Paolo Alfaro Sotil \\
            \textbf{Rol:} & Full-stack Developer \\
            \textbf{GitHub:} & \href{https://github.com/elbrujo325}{@elbrujo325} \\
            \textbf{LinkedIn:} & \href{https://linkedin.com/in/henry-paolo-alfaro-sotil}{linkedin.com/in/henry-paolo-alfaro-sotil} \\
            \textbf{Email:} & paolo.alfaro.sotil@gmail.com \\
            \textbf{Fecha:} & \DTMtoday \\
            \textbf{Versión:} & 1.0 \\
            \textbf{Repositorio:} & \href{https://github.com/elbrujo325/landing-maria-fernanda}{github.com/elbrujo325/landing-maria-fernanda} \\
            \textbf{Deploy:} & \href{https://landing-maria-fernanda.web.app}{landing-maria-fernanda.web.app}
        \end{tabular}
    \end{minipage}

    \vfill

    {\color{textmuted}\textit{Desarrollado con enfoque en performance, accesibilidad y mantenibilidad.}}

    \vspace{1cm}
\end{titlepage}

% Tabla de contenidos
\tableofcontents
\clearpage

% Lista de figuras y tablas
\listoffigures
\listoftables
\clearpage

\chapter{Introducción y Visión General}
\label{ch:introduccion}

\section{Propósito del Documento}
Este documento proporciona una documentación técnica exhaustiva del proyecto \textbf{Landing María Fernanda}, una página de aterrizaje (landing page) profesional para la psicóloga clínica María Fernanda Arana, especialista en Terapia de Aceptación y Compromiso (ACT).

Está dirigido a dos audiencias:
\begin{itemize}
    \item \textbf{Cliente (María Fernanda Arana):} Para comprender el alcance, funcionalidades y cómo gestionar el contenido.
    \item \textbf{Desarrolladores:} Para onboarding rápido, mantenimiento, escalabilidad y toma de decisiones técnicas.
\end{itemize}

\section{Alcance del Proyecto}
\begin{table}[H]
    \centering
    \caption{Alcance funcional del proyecto}
    \label{tab:alcance}
    \begin{tabular}{>{\bfseries}ll}
        \toprule
        \textbf{Área} & \textbf{Descripción} \\
        \midrule
        Landing Principal & Hero, Sobre mí, Credenciales, Problema/Solución, Servicios, Testimonios, FAQ, Footer \\
        Blog Público & Listado con filtros, búsqueda, paginación infinita, vista detalle con posts relacionados \\
        Panel Admin & Login seguro, CRUD posts, CRUD testimonios, editor Quill.js rich-text \\
        Backend & Firebase (Firestore, Auth, Hosting) — Plan Spark gratuito \\
        Performance & WebP srcset (8 densidades), lazy loading, image-rendering crisp, zero-runtime JS \\
        Accesibilidad & WCAG 2.1 AA, semántica HTML5, ARIA labels, focus visible, contraste \\
    \bottomrule
    \end{tabular}
\end{table}

\section{Objetivos de Negocio}
\begin{enumerate}
    \item \textbf{Conversión principal:} Agendar sesión vía WhatsApp (float button + sticky CTA móvil + hero buttons)
    \item \textbf{Autoridad clínica:} Mostrar credenciales, formación ACT, experiencia (+7 años, +12 países)
    \item \textbf{Content Marketing:} Blog SEO-ready para atraer tráfico orgánico (ansiedad, relaciones, propósito)
    \item \textbf{Autonomía:} Panel admin para que la clienta gestione posts/testimonios sin dependencia técnica
\end{enumerate}

\section{Stack Tecnológico Resumido}
\begin{table}[H]
    \centering
    \caption{Stack tecnológico}
    \label{tab:stack}
    \begin{tabular}{lll}
        \toprule
        \textbf{Capa} & \textbf{Tecnología} & \textbf{Justificación} \\
        \midrule
        Frontend & HTML5, CSS3 (Custom Props), Vanilla JS (ESM) & Cero dependencias, performance máximo, control total \\
        BaaS & Firebase (Firestore, Auth, Hosting) & Serverless, escalable, plan gratuito generoso, DX excelente \\
        Editor & Quill.js v2 (CDN) & Rich-text probado, modular, image-resize nativo \\
        Imágenes & WebP + srcset (72--800px), square-crop & Core Web Vitals (LCP, CLS), retina-ready \\
        Despliegue & Firebase CLI + GitHub (SSH) & CI/CD simple, preview channels, rollback instantáneo \\
    \bottomrule
    \end{tabular}
\end{table}

\chapter{Arquitectura del Sistema}
\label{ch:arquitectura}

\section{Diagrama de Arquitectura}
\begin{figure}[H]
    \centering
    \fbox{\begin{minipage}{0.9\textwidth}
        \centering
        \texttt{
        +---------------------------------------------------------------------+\textbackslash\\
        |                        USUARIO FINAL                                |\textbackslash\\
        |  (Navegador: Chrome, Firefox, Safari, Edge -- Mobile/Desktop)      |\textbackslash\\
        +---------------------------------------------------------------------+\textbackslash\\
                                  | HTTPS\textbackslash\\
                                  v\textbackslash\\
        +---------------------------------------------------------------------+\textbackslash\\
        |                  FIREBASE HOSTING (CDN Global)                      |\textbackslash\\
        |  public/index.html, blog.html, blog-post.html, admin*.html          |\textbackslash\\
        |  assets/images/*.webp, assets/logos/*, src/js/firebase-config.js    |\textbackslash\\
        +---------------------------------------------------------------------+\textbackslash\\
                                  |\textbackslash\\
          +-----------------------+-----------------------+\textbackslash\\
          v                       v                       v\textbackslash\\
    +-------------+       +-------------+       +-------------+\textbackslash\\
    | FIRESTORE   |       | FIREBASE    |       |  STORAGE    |\textbackslash\\
    | (NoSQL DB)  |       |   AUTH      |       | (Opcional)  |\textbackslash\\
    |             |       |             |       |             |\textbackslash\\
    | blogPosts   |<----->| Email/      |       | Imagenes    |\textbackslash\\
    | testimonios | Admin | Password    |       | Quill       |\textbackslash\\
    | (public/    |claims | + Custom    |       | (base64/    |\textbackslash\\
    |  admin)     |admin: |  claims     |       |  URL ext)   |\textbackslash\\
    +-------------+ true  +-------------+       +-------------+\textbackslash\\
        }
    \end{minipage}}
    \caption{Arquitectura de alto nivel -- Landing María Fernanda}
    \label{fig:arquitectura}
\end{figure}

\section{Modelo de Datos (Firestore)}

\subsection{Colección: \texttt{blogPosts}}
\begin{table}[H]
    \centering
    \caption{Campos de documento blogPosts}
    \label{tab:blogposts}
    \begin{tabular}{lllp{5cm}}
        \toprule
        \textbf{Campo} & \textbf{Tipo} & \textbf{Req.} & \textbf{Descripción} \\
        \midrule
        titulo & string & Sí & Título del post (SEO, max 80 chars) \\
        slug & string & Sí & URL amigable, único, minúsculas, guiones \\
        extracto & string & Sí & Resumen para cards/listado (max 200 chars) \\
        contenido & string & Sí & HTML completo desde Quill.js \\
        categoria & string & Sí & Enum: ansiedad, relaciones, autoestima, proposito, tecnicas, general \\
        etiquetas & array<string> & No & Tags para filtrado/búsqueda \\
        imagenDestacada & string & No & URL WebP (hero del post) \\
        publicado & boolean & Sí & \texttt{true} = visible en blog.html/blog-post.html \\
        fecha & timestamp & Sí & Fecha publicación (orden descendente) \\
        fechaCreacion & timestamp & Sí & Auditoría \\
        fechaActualizacion & timestamp & No & Auditoría \\
        autorUid & string & Sí & FK a Firebase Auth UID \\
        tiempoLectura & number & No & Minutos estimados (cálculo automático) \\
        \bottomrule
    \end{tabular}
\end{table}

\subsection{Colección: \texttt{testimonios}}
\begin{table}[H]
    \centering
    \caption{Campos de documento testimonios}
    \label{tab:testimonios}
    \begin{tabular}{lllp{5cm}}
        \toprule
        \textbf{Campo} & \textbf{Tipo} & \textbf{Req.} & \textbf{Descripción} \\
        \midrule
        nombre & string & Sí & Nombre o iniciales del paciente \\
        texto & string & Sí & Testimonio completo \\
        avatar & string & No & URL imagen (SVG placeholder si ausencia) \\
        iniciales & string & Sí & 2-3 letras para avatar fallback \\
        publicado & boolean & Sí & Visible en index.html \#testimonios \\
        orden & number & Sí & Orden de aparición (0, 1, 2...) \\
        fechaCreacion & timestamp & Sí & Auditoría \\
        \bottomrule
    \end{tabular}
\end{table}

\section{Índices Compuestos Requeridos}
Definidos en \texttt{firestore.indexes.json}:

\begin{lstlisting}[language=JSON, caption={firestore.indexes.json}]
{
  "indexes": [
    {
      "collectionGroup": "blogPosts",
      "queryScope": "COLLECTION",
      "fields": [
        { "fieldPath": "publicado", "order": "ASCENDING" },
        { "fieldPath": "fecha", "order": "DESCENDING" }
      ]
    },
    {
      "collectionGroup": "testimonios",
      "queryScope": "COLLECTION",
      "fields": [
        { "fieldPath": "publicado", "order": "ASCENDING" },
        { "fieldPath": "orden", "order": "ASCENDING" }
      ]
    }
  ]
}
\end{lstlisting}

\textbf{Importante:} Sin el índice \texttt{blogPosts (publicado Asc, fecha Desc)}, la consulta en \texttt{blog.html} fallará con error \texttt{FAILED\_PRECONDITION}. Crear desde Firebase Console o \texttt{firebase deploy --only firestore:indexes}.

\section{Reglas de Seguridad (Firestore)}
\begin{lstlisting}[language=JavaScript, caption={firestore.rules (resumen)}]
rules_version = '2';
service cloud.firestore {
  match /databases/{database}/documents {
    // Blog Posts
    match /blogPosts/{docId} {
      // Público: solo lee publicados
      allow read: if resource.data.publicado == true || request.auth != null;
      // Admin: escribe todo (custom claim admin:true)
      allow write: if request.auth != null && request.auth.token.admin == true;
    }

    // Testimonios
    match /testimonios/{docId} {
      allow read: if resource.data.publicado == true || request.auth != null;
      allow write: if request.auth != null && request.auth.token.admin == true;
    }
  }
}
\end{lstlisting}

\chapter{Frontend — Estructura y Componentes}
\label{ch:frontend}

\section{Páginas HTML — Responsabilidades}
\begin{table}[H]
    \centering
    \caption{Páginas y su propósito}
    \label{tab:paginas}
    \begin{tabular}{llp{6cm}}
        \toprule
        \textbf{Archivo} & **Rol** & \textbf{Descripción} \\
        \midrule
        index.html & Landing Principal & Hero, Sobre mí, Credenciales, Problema/Solución, Servicios, Testimonios, FAQ, Footer Luxto \\
        blog.html & Listado Blog & Grid responsivo, filtros categoría, búsqueda temps real, scroll infinito, footer Luxto \\
        blog-post.html & Detalle Post & Artículo completo, barra progreso lectura, share, posts relacionados (3), CTA final, footer Luxto \\
        admin.html & Dashboard Admin & Login, tabs Posts/Testimonios, CRUD modal, stats, footer Luxto \\
        admin-editor.html & Editor Quill & Toolbar completa, image resize, clipboard cleanup, preview, auto-save, footer Luxto \\
        \bottomrule
    \end{tabular}
\end{table}

\section{Sistema de Diseño (Design Tokens)}
Variables CSS en \texttt{:root} (\texttt{index.html:1-80}):

\begin{lstlisting}[language=CSS, caption={Design Tokens — Colores y Espaciado}]
:root {
  /* Colores principales */
  --celeste-deep: #11698e;      /* Primary brand */
  --celeste-mid:  #3aa8c9;
  --celeste-light: #6fc1ff;     /* Accent / hover */
  --cream:        #fff8e7;      /* Background warm */
  --cream-light:  #fffef0;
  --yellow-mid:   #ffe066;      /* Highlights */
  --white:        #ffffff;
  --text-dark:    #1a1a2e;
  --text-muted:   #6b7280;
  --card-border:  rgba(17,105,142,0.12);
  --card-bg:      #ffffff;

  /* Espaciado (escala 4px) */
  --space-xs: 4px;  --space-sm: 8px;  --space-md: 16px;
  --space-lg: 24px; --space-xl: 32px; --space-2xl: 48px;
  --space-3xl: 64px; --space-4xl: 96px;

  /* Tipografía */
  --font-heading: 'League Spartan', sans-serif;
  --font-body:    'Open Sans', sans-serif;

  /* Radii & Sombras */
  --radius-sm: 8px;  --radius-md: 12px; --radius-lg: 16px;
  --radius-xl: 24px; --radius-full: 9999px;
  --shadow-sm: 0 1px 2px rgba(0,0,0,0.05);
  --shadow-md: 0 4px 12px rgba(17,105,142,0.10);
  --shadow-lg: 0 12px 32px rgba(17,105,142,0.15);

  /* Colores sociales (footer) */
  --color-wa: #25d366;   --color-ig: #e1306c;
  --color-tt: #000000;   --color-fb: #1877f2;
}
\end{lstlisting}

\section{Breakpoints Responsive}
\begin{table}[H]
    \centering
    \caption{Breakpoints y comportamientos}
    \label{tab:breakpoints}
    \begin{tabular}{llp{6cm}}
        \toprule
        \textbf{Breakpoint} & \textbf{Dispositivo} & \textbf{Cambios Clave} \\
        \midrule
        $\geq$1440px & Desktop Large & Container 1280px, hero 2-col, footer 3-col (1.5/1/1fr) \\
        1024--1439px & Desktop/Tablet Landscape & Hero 2-col, footer 3-col \\
        768--1023px & Tablet Portrait & \textbf{Navbar $\rightarrow$ Hamburger absolute dropdown}, footer 2-col, hero 1-col \\
        $<$768px & Mobile & Footer 1-col, hero imagen arriba, CTA sticky, WhatsApp float push-up \\
        \bottomrule
    \end{tabular}
\end{table}

\section{Navbar — Comportamiento Detallado}
\begin{itemize}
    \item \textbf{Desktop ($>$768px):} Links horizontales, dropdown "Servicios" en hover/focus, logo + CTA "Agendar"
    \item \textbf{Mobile ($\leq$768px):} Hamburger ($\equiv$) $\rightarrow$ dropdown absoluto (\texttt{top: 100\%, left: 0, right: 0, z-index: 1000})
    \item \textbf{Animación:} \texttt{slideDown 0.3s cubic-bezier(0.4, 0, 0.2, 1)}
    \item \textbf{Scroll effect:} Clase \texttt{.scrolled} añade \texttt{backdrop-filter: blur(12px)} + sombra
    \item \textbf{Active link:} \texttt{IntersectionObserver} en secciones \texttt{id} resalta link actual
    \item \textbf{Bug fix crítico:} Selector "close menu on link click" excluye \texttt{.has-dropdown > a} para permitir abrir submenú
\end{itemize}

\section{Footer — Luxto Style (Referencia: index.html)}
Estructura de 3 zonas:

\begin{enumerate}
    \item \textbf{Franja Cita Centrada:} \texttt{blockquote} con frase "Si mejoras tú, mejora el mundo y los que te rodean" — \textit{María Fernanda Arana}
    \item \textbf{Grid 3 Columnas (1.5fr / 1fr / 1fr):}
    \begin{itemize}
        \item \textbf{Col 1 - Brand:} Logo círculo blanco + "Pensar Sentir Hacer" + Tagline "PSICOLOGÍA CLÍNICA ACT" + Descripción + Pill "Terapia con propósito"
        \item \textbf{Col 2 - Nav:} 7 links (Sobre mí, Servicios, ¿Te sientes así?, Testimonios, Blog, Preguntas, Agendar)
        \item \textbf{Col 3 - Contacto:} Ubicación (Lima, Perú), WhatsApp clickable, Email mailto, Horario + 4 Social Buttons (IG, TikTok, FB, WA) con \textbf{brand colors en hover}
    \end{itemize}
    \item \textbf{Bottom Bar:} Copyright © 2026 + Legal links (Privacidad, Términos, Behold)
\end{enumerate}

\textbf{Textura de ruido:} SVG filter \texttt{feTurbulence} en \texttt{.footer-background::before} (opacity 0.03).

\textbf{Responsive Footer:}
\begin{itemize}
    \item 1023px: Grid 2-col (Brand 100% width, Nav + Contacto 50% each)
    \item 767px: Grid 1-col, orden: Brand $\rightarrow$ Nav $\rightarrow$ Contacto
\end{itemize}

\section{WhatsApp Float — Colisión Inteligente}
\begin{lstlisting}[language=JavaScript, caption={WhatsApp Float collision detection}]
function updateWhatsAppFloat() {
  const whatsappFloat = document.querySelector('.whatsapp-float');
  const footer = document.querySelector('.footer');
  if (!whatsappFloat || !footer) return;

  const floatRect = whatsappFloat.getBoundingClientRect();
  const footerRect = footer.getBoundingClientRect();
  const viewportHeight = window.innerHeight;

  // Si el float toca o pasa el footer (con margen 16px)
  if (floatRect.bottom >= footerRect.top - 16) {
    const overlap = floatRect.bottom - footerRect.top + 16;
    whatsappFloat.style.bottom = `${overlap + 24}px`; // Push up
  } else {
    whatsappFloat.style.bottom = '24px'; // Posición normal
  }
}
// Throttled en scroll + resize
\end{lstlisting}

\chapter{Blog System — Implementación}
\label{ch:blog}

\section{Flujo de Datos}
\begin{figure}[H]
    \centering
    \fbox{\begin{minipage}{0.9\textwidth}
        \centering
        \texttt{
        ADMIN (admin-editor.html)                    PUBLICO (blog.html, blog-post.html)\textbackslash\\
        +---------------------+                      +---------------------+\textbackslash\\
        | 1. Escribe en Quill |                      | 1. onSnapshot()     |\textbackslash\\
        | 2. Click "Publicar" |---> Firestore ------>|    where(publicado  |\textbackslash\\
        |    ->               |   blogPosts          |    == true)         |\textbackslash\\
        |    {publicado:true} |   (trigger)          |    orderBy(fecha    |\textbackslash\\
        | 3. Doc creado/actual|                      |    desc)            |\textbackslash\\
        +---------------------+                      | 2. Render cards/    |\textbackslash\\
                                                     |    article HTML     |\textbackslash\\
                                                     +---------------------+\textbackslash\\
        }
    \end{minipage}}
    \caption{Flujo de datos Blog: Admin $\rightarrow$ Firestore $\rightarrow$ Publico}
    \label{fig:blog-flow}
\end{figure}

\section{blog.html — Listado}
\begin{itemize}
    \item \textbf{Query:} \texttt{db.collection('blogPosts').where('publicado', '==', true).orderBy('fecha', 'desc')}
    \item \textbf{Listener real-time:} \texttt{onSnapshot()} — actualización instantánea sin refresh
    \item \textbf{Filtros:} Chips por categoría (Todos, Ansiedad, Relaciones, Autoestima, Propósito, Técnicas, General)
    \item \textbf{Búsqueda:} Input \texttt{type="search"} con \texttt{debounce(300ms)} filtra por título/extracto/etiquetas
    \item \textbf{Paginación infinita:} \texttt{IntersectionObserver} en sentinel $\rightarrow$ \texttt{limit(6)} + \texttt{startAfter(lastDoc)}
    \item \textbf{Empty state:} Ilustración SVG + mensaje "No hay posts publicados" + CTA a admin (si logueado)
\end{itemize}

\section{blog-post.html — Detalle}
\begin{itemize}
    \item \textbf{Routing:} \texttt{blog-post.html?id=<docId>} (URL param)
    \item \textbf{Loading:} Spinner skeleton mientras \texttt{docSnap} carga
    \item \textbf{Not found:} 404 UI amigable + link "Volver al blog"
    \item \textbf{Reading Progress Bar:} Fixed top, width = \texttt{(scrollY / (articleHeight - viewportHeight)) * 100\%}
    \item \textbf{Share Buttons:} WhatsApp, LinkedIn, Twitter (X), Copy Link — \texttt{navigator.share()} fallback
    \item \textbf{Related Posts:} 3 posts misma categoría (excluyendo actual), fallback aleatorios
    \item \textbf{CTA Final:} "¿Te resonó este artículo?" + WhatsApp button
    \item \textbf{SEO Meta:} \texttt{title}, \texttt{description}, \texttt{og:*}, \texttt{twitter:*}, \texttt{article:*} inyectados vía JS
\end{itemize}

\section{Admin Panel — CRUD}

\subsection{Autenticación}
\begin{itemize}
    \item Firebase Auth: Email/Password
    \item \textbf{Custom Claim requerido:} \texttt{admin: true} (set via Admin SDK)
    \item \texttt{onAuthStateChanged} $\rightarrow$ redirige a login si no autenticado o sin claim
    \item Logout: \texttt{auth.signOut()} + redirect a \texttt{admin.html}
\end{itemize}

\subsection{Editor Quill.js (admin-editor.html)}
\begin{lstlisting}[language=JavaScript, caption={Configuración Quill.js}]
import Quill from 'https://cdn.jsdelivr.net/npm/quill@2.0.0-dev.4/+esm';
import ImageResize from 'quill-image-resize-module';

Quill.register('modules/imageResize', ImageResize);

const quill = new Quill('#editor', {
  theme: 'snow',
  modules: {
    toolbar: [
      [{ header: [1, 2, 3, false] }],
      ['bold', 'italic', 'underline', 'strike'],
      [{ color: [] }, { background: [] }],
      [{ list: 'ordered' }, { list: 'bullet' }, { indent: '-1' }, { indent: '+1' }],
      ['link', 'image', 'video', 'blockquote', 'code-block'],
      ['clean']
    ],
    clipboard: {
      matchVisual: false // Limpia estilos Word/Google Docs
    },
    imageResize: {
      displayStyles: { 'background-color': 'transparent', 'border': 'none' },
      modules: ['Resize', 'DisplaySize', 'Toolbar']
    },
    history: { delay: 1000, maxStack: 100, userOnly: true }
  },
  placeholder: 'Escribe tu artículo aquí...'
});
\end{lstlisting}

\textbf{Imágenes en Quill:} Actualmente \texttt{base64} o URL externa. \textit{TODO: Upload a Firebase Storage + CDN}.

\chapter{Optimización de Imágenes}
\label{ch:imagenes}

\section{Estrategia Dual: Original vs Optimizado}

\subsection{index.html — Fotos Originales (Máxima Calidad)}
\begin{table}[H]
    \centering
    \caption{Uso de foto original en index.html}
    \label{tab:foto-original}
    \begin{tabular}{llll}
        \toprule
        \textbf{Sección} & \textbf{Selector} & \textbf{Tamaño CSS} & \textbf{Archivo} \\
        \midrule
        Hero & \texttt{.hero-photo img} & 400px círculo (desktop) & \texttt{foto-fernanda.png} (377$\times$661) \\
        Perfil & \texttt{.perfil-photo img} & 260px círculo sticky & \texttt{foto-fernanda.png} (377$\times$661) \\
        \bottomrule
    \end{tabular}
\end{table}
\textit{Decisión de cliente: mantener retrato original sin recorte ni compresión en zonas hero.}

\subsection{Blog / Blog-Post — Avatares Optimizados (Square Crop + WebP)}

\subsubsection{Proceso de Optimización}
\begin{enumerate}
    \item \textbf{Source:} \texttt{foto-fernanda.png} (377$\times$661, portrait, 2.1 MB)
    \item \textbf{Crop:} Square 1:1 desde \textbf{top-center} (cara centrada) $\rightarrow$ \texttt{foto-fernanda-square.png} (400$\times$400)
    \item \textbf{WebP densities:} 72, 108, 144, 216, 260, 400, 520, 800px (8 variantes)
    \item \textbf{Naming:} \texttt{foto-fernanda-\{w\}.webp}
\end{enumerate}

\subsubsection{HTML Implementación}
\begin{lstlisting}[language=HTML, caption={Avatar optimizado — blog.html / blog-post.html}]
<img src="assets/images/foto-fernanda-square.png"
     srcset="
       assets/images/foto-fernanda-72.webp   72w,
       assets/images/foto-fernanda-108.webp  108w,
       assets/images/foto-fernanda-144.webp  144w,
       assets/images/foto-fernanda-216.webp  216w,
       assets/images/foto-fernanda-260.webp  260w,
       assets/images/foto-fernanda-400.webp  400w,
       assets/images/foto-fernanda-520.webp  520w,
       assets/images/foto-fernanda-800.webp  800w
     "
     sizes="(max-width: 767px) 60px, (max-width: 1023px) 72px, 80px"
     alt="María Fernanda Arana"
     loading="lazy"
     style="image-rendering: -webkit-optimize-contrast; image-rendering: crisp-edges;"
     width="80" height="80">
\end{lstlisting}

\subsubsection{CSS para Nitidez en Pequeño}
\begin{lstlisting}[language=CSS, caption={image-rendering para avatares pequeños}]
.avatar-img,
.blog-author-avatar,
.post-author-avatar {
  image-rendering: -webkit-optimize-contrast; /* Safari/Chrome */
  image-rendering: crisp-edges;               /* Firefox/Edge */
  image-rendering: pixelated;                 /* Fallback */
}
\end{lstlisting}

\subsubsection{Resultado de Performance}
\begin{table}[H]
    \centering
    \caption{Métricas de imagen avatar (WebP vs PNG original)}
    \label{tab:img-perf}
    \begin{tabular}{lrrr}
        \toprule
        \textbf{Métrica} & \textbf{Original PNG} & \textbf{WebP 80px} & \textbf{Mejora} \\
        \midrule
        Tamaño archivo & 2.1 MB & 2.8 KB (72w) -- 18 KB (800w) & \textbf{99\% reducción} \\
        Dimensiones & 377$\times$661 & 72--800px (responsive) & \textbf{Srcset nativo} \\
        Formato & PNG (lossless) & WebP (lossy, quality 80) & \textbf{Modern browser support} \\
        Renderizado & Blurry en 60px & Crisp (image-rendering) & \textbf{Legibilidad avatar} \\
        \bottomrule
    \end{tabular}
\end{table}

\chapter{Despliegue y Operaciones}
\label{ch:despliegue}

\section{Comandos Firebase CLI}
\begin{lstlisting}[language=bash, caption={Comandos de despliegue}]
# Login (una vez)
firebase login

# Desarrollo local (emuladores)
firebase emulators:start --only hosting,firestore,auth

# Servir solo hosting (rápido)
firebase serve --only hosting --port 5000

# Deploy completo (Hosting + Firestore rules + indexes)
firebase deploy --project landing-maria-fernanda

# Deploy solo hosting
firebase deploy --only hosting --project landing-maria-fernanda

# Deploy solo reglas e índices
firebase deploy --only firestore:rules,firestore:indexes --project landing-maria-fernanda

# Ver logs hosting
firebase hosting:channel:list
firebase hosting:channel:open <channel-id>
\end{lstlisting}

\section{Configuración firebase.json}
\begin{lstlisting}[language=JSON, caption={firebase.json}]
{
  "hosting": {
    "public": "public",
    "ignore": [
      "firebase.json",
      "**/.*",
      "**/node_modules/**"
    ],
    "rewrites": [
      {
        "source": "**",
        "destination": "/index.html"
      }
    ],
    "headers": [
      {
        "source": "**/*.@(webp|avif|png|jpg|jpeg|svg|ico)",
        "headers": [
          {
            "key": "Cache-Control",
            "value": "public, max-age=31536000, immutable"
          }
        ]
      },
      {
        "source": "**/*.html",
        "headers": [
          {
            "key": "Cache-Control",
            "value": "public, max-age=0, must-revalidate"
          }
        ]
      }
    ]
  },
  "firestore": {
    "rules": "firestore.rules",
    "indexes": "firestore.indexes.json"
  },
  "emulators": {
    "auth": { "port": 9099 },
    "firestore": { "port": 8080 },
    "hosting": { "port": 5000 },
    "ui": { "enabled: true, "port": 4000 }
  }
}
\end{lstlisting}

\section{Git Workflow}
\begin{lstlisting}[language=bash, caption={Git workflow estándar}]
# Trabajo en feature branch
git checkout -b feat/nueva-funcionalidad
# ... cambios ...
git add .
git commit -m "feat: descripción clara del cambio"
git push origin feat/nueva-funcionalidad
# PR en GitHub $\rightarrow$ review $\rightarrow$ merge a main

# Deploy automático tras merge a main (opcional: GitHub Actions)
# O deploy manual:
git checkout main
git pull origin main
firebase deploy --project landing-maria-fernanda
\end{lstlisting}

\textbf{Remote SSH:} \texttt{git@github.com:elbrujo325/landing-maria-fernanda.git}

\section{Configuración Admin (Custom Claims)}
\begin{lstlisting}[language=bash, caption={Asignar claim admin via Firebase CLI}]
# Opción A: Firebase Admin SDK (Node)
# node setAdminClaim.js <uid>
# const admin = require('firebase-admin');
# admin.auth().setCustomUserClaims(uid, { admin: true });

# Opción B: Firebase Console
# 1. Authentication > Users > seleccionar usuario
# 2. Custom claims (requiere Blaze plan) o Cloud Function

# Opción C: Cloud Function (recomendado para producción)
# exports.setAdminClaim = functions.https.onCall(async (data, context) => {
#   if (!context.auth.token.admin) throw new Error('No autorizado');
#   await admin.auth().setCustomUserClaims(data.uid, { admin: true });
#   return { success: true };
# });
\end{lstlisting}

\chapter{Mantenimiento y Troubleshooting}
\label{ch:mantenimiento}

\section{Checklist Pre-Deploy}
\begin{table}[H]
    \centering
    \caption{Checklist de calidad antes de deploy}
    \label{tab:checklist}
    \begin{tabular}{lll}
        \toprule
        \textbf{Categoría} & \textbf{Verificación} & \textbf{Herramienta} \\
        \midrule
        HTML & Válido, semántico & \texttt{npx html-validate public/*.html} \\
        CSS & Sin errores, variables usadas & \texttt{npx stylelint public/**/*.html} \\
        JS & Sin errores consola & Chrome DevTools (todas las páginas) \\
        Responsive & 320, 768, 1024, 1440px & Chrome Device Toolbar \\
        Firestore & Índices verdes & Firebase Console > Firestore > Indexes \\
        Imágenes & WebP servido & Network tab > img > type: webp \\
        Accesibilidad & 0 critical/serious & axe DevTools / Lighthouse \\
        Performance & Lighthouse $\geq$ 90 & Lighthouse CI / DevTools \\
        SEO & Meta tags completos & Lighthouse SEO / Social Preview \\
        \bottomrule
    \end{tabular}
\end{table}

\section{Problemas Conocidos y Soluciones}

\begin{table}[H]
    \centering
    \caption{Troubleshooting común}
    \label{tab:troubleshooting}
    \begin{tabular}{p{3cm}p{5cm}p{5cm}}
        \toprule
        \textbf{Síntoma} & \textbf{Causa} & \textbf{Solución} \\
        \midrule
        Blog no carga posts & Índice compuesto faltante & Firebase Console > Firestore > Indexes > "Create index" para \texttt{blogPosts (publicado Asc, fecha Desc)} \\
        Admin: "Permission denied" & Falta custom claim \texttt{admin:true} & Asignar claim via Admin SDK o Cloud Function \\
        Quill: imagen no sube & No hay Storage configurado & Implementar upload a Firebase Storage + retornar URL \\
        Footer se rompe en móvil & \texttt{.container} wrapper en \texttt{.footer-grid} & Remover \texttt{.container} del footer HTML (CSS espera \texttt{.footer-grid} directo) \\
        Navbar móvil no abre & CSS breakpoint 968px vs JS 768px & Unificar a 768px en CSS y JS \\
        Dropdown "Servicios" no abre & Link capturado por "close on link click" & Excluir \texttt{.has-dropdown > a} del handler \\
        Avatar pixelado en blog & Usaba foto original 377x661 en circle 60px & Square crop top-center + WebP srcset + image-rendering:crisp-edges \\
        \bottomrule
    \end{tabular}
\end{table}

\section{Logs y Debugging}
\begin{itemize}
    \item \textbf{Firebase Hosting Logs:} \texttt{firebase hosting:channel:open <id>} o Console > Hosting > Logs
    \item \textbf{Firestore:} Console > Logs > Firestore (reads/writes, errors)
    \item \textbf{Client-side:} \texttt{localStorage.setItem('debug', 'true')} $\rightarrow$ logs verbose en consola
    \item \textbf{Performance:} \texttt{performance.mark() / measure()} en secciones críticas (hero, blog load)
\end{itemize}

\chapter{Escalabilidad y Roadmap}
\label{ch:roadmap}

## Deuda Técnica y Mejoras Priorizadas

\begin{table}[H]
    \centering
    \caption{Roadmap técnico}
    \label{tab:roadmap}
    \begin{tabular}{llll}
        \toprule
        \textbf{Prioridad} & \textbf{Item} & \textbf{Esfuerzo} & \textbf{Impacto} \\
        \midrule
        🔴 Crítico & Índice compuesto blogPosts en producción & 5 min & Blog funcional \\
        🟡 Alto & Quill image upload $\rightarrow$ Firebase Storage & 4h & Editor completo, sin base64 \\
        🟡 Alto & SEO meta tags dinámicos por post (SSR/SSG) & 8h & SEO orgánico real \\
        🟡 Medio & PWA: Manifest + Service Worker (offline) & 6h & UX móvil, engagement \\
        🟢 Bajo & Tests E2E (Playwright) & 16h & Confianza deploy \\
        🟢 Bajo & Analytics (GA4 + eventos custom) & 4h & Métricas negocio \\
        🟢 Bajo & i18n (EN/PT) para pacientes internacionales & 12h & Expansión mercado \\
        \bottomrule
    \end{tabular}
\end{table}

\section{Arquitectura Futura (v2.0)}
\begin{itemize}
    \item \textbf{Next.js + Firebase (App Router):} SSR/ISR para blog-post.html $\rightarrow$ SEO real, OG tags estáticos
    \item \textbf{Cloud Functions v2:} Triggers para \texttt{onCreate} blogPosts $\rightarrow$ generar OG images dinámicas, sitemap.xml, push notifications
    \item \textbf{Firebase Extensions:} \texttt{Resize Images} para Quill uploads automático
    \item \textbf{Algolia/Meilisearch:} Búsqueda full-text real en blog (reemplazar filtro cliente)
    \item \textbf{Headless CMS opcional:} Si cliente necesita multi-autor, workflow editorial
\end{itemize}

\chapter{Apéndices}
\label{ch:apendices}

\section{Archivos de Configuración Clave}

\subsection{package.json (si se usa build tools)}
\begin{lstlisting}[language=JSON, caption={package.json sugerido para tooling}]
{
  "name": "landing-maria-fernanda",
  "version": "1.0.0",
  "private": true,
  "scripts": {
    "dev": "firebase serve --only hosting",
    "deploy": "firebase deploy --project landing-maria-fernanda",
    "deploy:rules": "firebase deploy --only firestore:rules,firestore:indexes",
    "lint:html": "html-validate public/*.html",
    "lint:css": "stylelint public/**/*.html",
    "test:e2e": "playwright test",
    "format": "prettier --write \"public/**/*.{html,css,js}\""
  },
  "devDependencies": {
    "firebase-tools": "^13.0.0",
    "html-validate": "^8.0.0",
    "stylelint": "^16.0.0",
    "prettier": "^3.0.0",
    "@playwright/test": "^1.40.0"
  }
}
\end{lstlisting}

\subsection{.gitignore}
\begin{lstlisting}[language=bash, caption={.gitignore}]
# Firebase
.firebase/
firebase-debug.log
firestore-debug.log

# Dependencias
node_modules/

# IDE
.vscode/
.idea/
*.swp

# OS
.DS_Store
Thumbs.db

# Logs
*.log
npm-debug.log*

# Testing
coverage/
.nyc_output/
playwright-report/
test-results/

# Build (si aplica)
dist/
build/
\end{lstlisting}

\section{Referencias y Recursos}

\begin{enumerate}
    \item \textbf{Firebase Docs:} \url{https://firebase.google.com/docs}
    \item \textbf{Firestore Data Modeling:} \url{https://firebase.google.com/docs/firestore/data-model}
    \item \textbf{Quill.js v2 Docs:} \url{https://quilljs.com/docs/}
    \item \textbf{WebP Browser Support:} \url{https://caniuse.com/webp}
    \item \textbf{CSS Custom Properties:} \url{https://developer.mozilla.org/en-US/docs/Web/CSS/--*}
    \item \textbf{IntersectionObserver API:} \url{https://developer.mozilla.org/en-US/docs/Web/API/IntersectionObserver_API}
    \item \textbf{WCAG 2.1 Guidelines:} \url{https://www.w3.org/WAI/WCAG21/quickref/}
    \item \textbf{Luxto Design Inspiration:} \url{https://luxto.com/} (footer style reference)
\end{enumerate}

\section{Contacto y Soporte}

Para dudas técnicas, mejoras o incidencias:

\begin{itemize}
    \item \textbf{Desarrollador:} Henry Paolo Alfaro Sotil
    \item \textbf{GitHub:} \href{https://github.com/elbrujo325}{@elbrujo325} — Issues y PRs bienvenidos
    \item \textbf{LinkedIn:} \href{https://linkedin.com/in/henry-paolo-alfaro-sotil}{linkedin.com/in/henry-paolo-alfaro-sotil}
    \item \textbf{Email:} paolo.alfaro.sotil@gmail.com
    \item \textbf{Repositorio:} \href{https://github.com/elbrujo325/landing-maria-fernanda}{github.com/elbrujo325/landing-maria-fernanda}
\end{itemize}

\vspace{2cm}

\begin{center}
    \rule{0.5\textwidth}{0.5pt}\\
    \vspace{0.5cm}
    \textit{Documentación generada el \DTMtoday \quad Versión 1.0}\\
    \textit{Landing María Fernanda — Psicología Clínica ACT}\\
    \textit{Desarrollado por Henry Paolo Alfaro Sotil}
\end{center}

\end{document}